\documentclass[11pt]{article}

\usepackage[margin=0.92in]{geometry}
\usepackage{amsmath,amssymb,amsthm,mathtools,booktabs,enumitem,microtype}
\usepackage{xcolor}
\usepackage[colorlinks=true,linkcolor=blue!50!black,urlcolor=blue!50!black,citecolor=blue!50!black]{hyperref}

\theoremstyle{plain}
\newtheorem{theorem}{Theorem}
\newtheorem{lemma}[theorem]{Lemma}
\newtheorem{proposition}[theorem]{Proposition}
\newtheorem{corollary}[theorem]{Corollary}

\theoremstyle{definition}
\newtheorem{definition}[theorem]{Definition}

\theoremstyle{remark}
\newtheorem{remark}{Remark}

\DeclareMathOperator{\Var}{Var}

\title{\bfseries The Sharp Cubic Rademacher Stability Constant}
\author{}
\date{\vspace{-1.6em}%
\small
\begin{tabular}{@{}c@{}}
Target paper: P.~Gao and J.~Qian, \emph{Fourth-Moment Geometry of Rademacher Sums},\\
\href{https://arxiv.org/abs/2608.17802}{arXiv:2608.17802} (August 2026), Conjecture~\texttt{eq:conjecture}\\[0.35em]
Companion verification script: \texttt{solutions/rademacher\_fourth\_moment/explore.py}\\[0.35em]
November 2026
\end{tabular}}

\begin{document}

\maketitle

\begin{abstract}
We prove Conjecture~\texttt{eq:conjecture} of Gao and Qian (\href{https://arxiv.org/abs/2608.17802}{arXiv:2608.17802}, August~2026) in full generality for all dimensions $n \ge 1$ with the sharp dimension-free constant $C = 5\sqrt{3} - 6\sqrt{2}$: for every unit vector $a \in S^{n-1}$ and i.i.d.\ uniform signs $\varepsilon_1, \dots, \varepsilon_n \in \{\pm 1\}$,
\[
\mathbb{E}\bigl|\overline{S}_n\bigr|^3 - \mathbb{E}\!\left|\sum_{i=1}^n a_i \varepsilon_i\right|^3
\;\ge\; (5\sqrt{3} - 6\sqrt{2})\sum_{i=1}^n \Bigl(a_i^2 - \frac{1}{n}\Bigr)^{\!2},
\]
with nonflat equality at $n = 3$, $a = (1, 1, 0)/\sqrt{2}$.
Our computer-assisted proof combines a scalar Fourier majorant (settling all even $n \ge 6$ and far odd $n \ge 7$ via $6$ exact rational endpoints), a two-coordinate equalization flow driven by a central-slice Jensen bound certified simultaneously for all odd $n \ge 7$ via $6{,}609$ integer-arithmetic interval boxes, and exact integer polynomial cone certificates ($47$ leaves total) yielding the dimension-specific optimal constants for $2 \le n \le 5$.
Every numerical bound in the proof relies solely on exact integer and rational arithmetic.
\end{abstract}

\section{Statement, Sharpness, and Exact Rational Preliminaries}

Let $\varepsilon_1, \dots, \varepsilon_n$ be i.i.d.\ uniform $\pm 1$ signs and $a \in S^{n-1}$.
Absorbing signs of $a_i$ into $\varepsilon_i$, set $x_i := a_i^2 \ge 0$ (so $\sum_{i=1}^n x_i = 1$), $q := \sum_{i=1}^n x_i^2$, $\Delta_n(a) := \sum_{i=1}^n (x_i - 1/n)^2 = q - 1/n$,
\[
F(x) := \mathbb{E}\!\left|\sum_{i=1}^n \sqrt{x_i}\,\varepsilon_i\right|^3, \qquad
M_n := \mathbb{E}\bigl|\overline{S}_n\bigr|^3, \qquad
\overline{S}_n := n^{-1/2}\sum_{i=1}^n \varepsilon_i.
\]
Let $C := 5\sqrt{3} - 6\sqrt{2}$ and define the rational threshold $c := 7/40$.
All terminating decimals in inequalities denote exact rational numbers.

\begin{theorem}[Sharp cubic Rademacher stability]\label{thm:main}
For every integer $n \ge 1$ and every $a \in S^{n-1}$,
\begin{equation}\label{eq:main}
M_n - F(x) \;\ge\; C\,\Delta_n(a), \qquad C = 5\sqrt{3} - 6\sqrt{2}.
\end{equation}
The constant $C$ is sharp: at $n = 3$ and $a = (1, 1, 0)/\sqrt{2}$, $M_3 = 5/(2\sqrt{3})$, $F(x) = \sqrt{2}$, and $\Delta_3(a) = 1/6$, giving $(M_3 - F(x))/\Delta_3(a) = 5\sqrt{3} - 6\sqrt{2}$.
\end{theorem}

\noindent\textbf{Exact rational bounds.}
Using $3.14159 < \pi < 3.14160$ and exact integer arithmetic, we record once and for all:
\begin{itemize}[leftmargin=1.5em,itemsep=2pt]
\item $\frac{4}{\pi}\bigl(\frac{0.64}{4} - \frac{1}{32}\bigr) > \frac{4}{3.14160}\cdot\frac{103}{800} = \frac{10300}{62832} > 0.16$;
\item $1.595^2 = 2.544025 < \frac{8}{3.14160} < \frac{8}{\pi} < \frac{8}{3.14159} < 2.547216 = 1.596^2$, hence $1.595 < \mu := \sqrt{8/\pi} < 1.596$;
\item $0.797^2 = 0.635209 < \frac{2}{3.14160} < \frac{2}{\pi}$, hence $\sqrt{2/\pi} > 0.797$;
\item $C < \frac{7}{40}$, since $5\sqrt{3} < 6\sqrt{2} + \frac{7}{40} \iff \frac{4751}{1600} < \frac{3360}{1600}\sqrt{2} \iff 4751^2 = 22{,}572{,}001 < 22{,}579{,}200 = 2 \cdot 3360^2$;
\item For the small-dimension constants in Section~\ref{sec:small}: $C_{2,1} = 2(\sqrt{2}-1) > \frac{7}{40}$ ($12{,}800 > 87^2 = 7{,}569$), $C_{4,2} = 6 - 4\sqrt{2} > \frac{7}{40}$ ($233^2 = 54{,}289 > 51{,}200$), and $C_{5,4} = \frac{27\sqrt{5}}{2} - 30 > \frac{7}{40}$ ($5 \cdot 540^2 = 1{,}458{,}000 > 1207^2 = 1{,}456{,}849$).
\end{itemize}

\section{A Scalar Fourier Majorant}

For $s > 0$ and $x > 0$, define
\[
p(s) := \frac{2}{\pi}\int_0^{\pi/2} \cos^s t\,dt, \qquad
h(s) := 2\sqrt{s}\,p(s), \qquad
g(x) := x\,h(1/x), \quad g(0) := 0.
\]

\begin{proposition}[Fourier majorant]\label{prop:fourier}
For every $x \in \Delta^{n-1}$,
\begin{equation}\label{eq:majorant}
F(x) \;\le\; \sum_{i=1}^n g(x_i), \qquad\text{and}\qquad M_n = h(n) \quad (\text{even } n).
\end{equation}
\end{proposition}

\begin{proof}
Set $J := \int_0^\infty (\cos t - 1 + t^2/2)t^{-4}\,dt \in (0, \infty)$.
Because $(\cos t - 1 + t^2/2)t^{-4} \ge 0$, scaling gives $J^{-1}\int_0^\infty (\cos(ut) - 1 + u^2 t^2/2)t^{-4}\,dt = |u|^3$ for all $u \in \mathbb{R}$, and Tonelli's theorem gives $F(x) = J^{-1}\int_0^\infty \bigl(\prod_{i=1}^n \cos(\sqrt{x_i}\,t) - 1 + t^2/2\bigr)t^{-4}\,dt$.
By weighted AM--GM over nonzero coordinates,
\[
\prod_{i=1}^n \cos(\sqrt{x_i}\,t) \;\le\; \prod_{i: x_i > 0} |\cos(\sqrt{x_i}\,t)| \;\le\; \sum_{i: x_i > 0} x_i |\cos(\sqrt{x_i}\,t)|^{1/x_i},
\]
with equality at $x_i = 1/n$ when $n$ is even (as $\cos^n(t/\sqrt{n}) = |\cos(t/\sqrt{n})|^n$).
Each term on the right has quadratic Taylor term $-t^2/2$; subtracting $1 - t^2/2$, integrating against $J^{-1}t^{-4}\,dt$, and rescaling reduce~\eqref{eq:majorant} to
\begin{equation}\label{eq:fourier-id}
J^{-1}\int_0^\infty \bigl(|\cos t|^s - 1 + \tfrac{1}{2}st^2\bigr)t^{-4}\,dt = 2s^2 p(s) \qquad (s > 0).
\end{equation}

For $s > 6$, $f(t) = |\cos t|^s \in C^6(\mathbb{R})$ is even and $\pi$-periodic, so its Fourier series $\sum_{j \in \mathbb{Z}} b_j e^{2ijt}$ has $b_{-j} = b_j$, $b_0 = p(s)$, and $b_j = O_s(|j|^{-6})$ by six integrations by parts.
From $f'\cos t = -s f\sin t$ we obtain $(u + j + 1)b_{j+1} = (u - j)b_j$ with $u := s/2$.
Putting $P(j) := \frac{1}{2}[j^3 + (u - 1)j^2 + u^2]$, direct expansion gives $(u + j + 1)P(j) - (u - j)P(j + 1) = (u + j + 1)j^3$, hence $j^3 b_j = P(j)b_j - P(j + 1)b_{j+1}$.
Telescoping yields $\sum_{j \ge 0} j^3 b_j = P(0)b_0 = \frac{1}{2}u^2 b_0$.
Since $f(0) = 1$ and $f''(0) = -s$, absolute convergence ($\sum |j|^3|b_j| < \infty$) and the nonnegative integration kernel justify termwise integration in~\eqref{eq:fourier-id}, giving $\sum_{j \in \mathbb{Z}} |2j|^3 b_j = 16 \cdot \frac{1}{2}u^2 b_0 = 2s^2 p(s)$.
Both sides of~\eqref{eq:fourier-id} are holomorphic on $\{\Re s > 0\}$: on compact subsets of $\{\Re s > 0\}$, the numerator is uniformly $O_K(t^4)$ near $0$ and dominated by $2 + \frac{1}{2}|s|t^2$ on $[1, \infty)$ (setting $|\cos t|^s := 0$ at zeros of $\cos t$).
Dominated convergence and the identity theorem extend~\eqref{eq:fourier-id} to all $s > 0$.
\end{proof}

\section{Curvature, Even Dimensions \texorpdfstring{$n \ge 6$}{n >= 6}, and Far Odd Dimensions}

Integration by parts gives $p(s + 2) = \frac{s + 1}{s + 2}p(s)$, and rescaling by $\sqrt{s}$ with $\cos t \le e^{-t^2/2}$ on $[0, \pi/2]$ gives $h(s) \to \mu = \sqrt{8/\pi}$ as $s \to \infty$.

\begin{lemma}[Curvature bounds]\label{lem:curv}
For every $s > 0$,
\begin{equation}\label{eq:logh}
(\log h)'(s) = \frac{1}{2}\int_0^\infty e^{-st}\tanh(t/2)\,dt.
\end{equation}
For $x = 1/s \in (0, 1]$, set $A_s := s\int_0^\infty u e^{-u}\tanh(u/(2s))\,du$ and $B_s := \int_0^\infty e^{-u}\tanh(u/(2s))\,du$. Then
\begin{equation}\label{eq:curv}
-\frac{g''(x)}{2} = h(s)\left(\frac{A_s}{4} - \frac{s B_s^2}{8}\right)
\;\ge\; \max\bigl\{0.16,\; L_0(x)\bigr\},
\end{equation}
where $L_0(x) := 1.595(1 - x/4)\bigl((1 - x^2)/4 - x/32\bigr)$.
\end{lemma}

\begin{proof}
Set $L(s) := \log\mu - \frac{1}{2}\int_0^\infty e^{-st}\tanh(t/2)\,dt/t$.
Then $(L(s) - L(s+2))' = \frac{1}{2}\int_0^\infty e^{-st}(1 - e^{-t})^2\,dt = \bigl(\log\frac{\sqrt{s(s+2)}}{s+1}\bigr)'$.
Since both $L(s) - L(s+2)$ and $\log(h(s)/h(s+2))$ vanish at infinity and share the same derivative, $L - \log h$ is $2$-periodic with limit~$0$, proving~\eqref{eq:logh}.

Differentiating~\eqref{eq:logh} and using $g''(x) = s^3 h''(s)$ gives the identity in~\eqref{eq:curv}.
Concavity of $\tanh$ on $[0, \infty)$ with $\tanh 0 = 0$ implies $A_s \ge A_1 = 1 - 2\sum_{j \ge 1}(-1)^{j-1}/(j+1)^2$.
The 25-term alternating upper bound for the sum yields $A_1 > 0.64$ (verified in exact rational arithmetic).
Together with $B_s \le 1/(2s)$ and $h(s) \ge h(1) = 4/\pi$, this gives $-g''(x)/2 \ge \frac{4}{\pi}(0.64/4 - 1/32) > 0.16$.
Next, $\tanh v \ge v - v^3/3$ gives $A_s \ge 1 - s^{-2}$, while~\eqref{eq:logh} and $\tanh v \le v$ give $h(s) \ge \mu e^{-1/(4s)} \ge 1.595(1 - 1/(4s))$, establishing $-g''(x)/2 \ge L_0(x)$ wherever the bracket is positive (and $0.16$ covers the rest).
\end{proof}

\begin{theorem}[Even $n \ge 6$ and far odd $n \ge 7$]\label{thm:far}
Inequality~\eqref{eq:main} holds with the stronger constant $c = 7/40 > C$ for:
\begin{enumerate}[label=\textup{(\arabic*)}]
\item every even $n \ge 6$ and all $x \in \Delta^{n-1}$;
\item every odd $n \ge 7$ whenever $D := n^2\Delta_n(a) = \sum_{i=1}^n (n x_i - 1)^2 \ge H(n)$, where $H(7) = 1.75$, $H(9) = 1.45$, and $H(n) = 1.24$ for odd $n \ge 11$.
\end{enumerate}
\end{theorem}

\begin{proof}
At $x_0 = 1/n$, the tangent-line deficit $g(x_0) + g'(x_0)(x - x_0) - g(x)$ equals $(x - x_0)^2$ times the average of $-g''/2$ along $[x_0, x]$ with weight $2(1 - t)$ ($t \in [0, 1]$), extending continuously to $x = 0$.
The first third $t \in [0, 1/3]$ carries weight $5/9$ and stays in $[0, \frac{2}{3n} + \frac{1}{3}]$.
Since $L_0$ decreases on $[0, 1/2]$, the tangent deficit coefficient is at least $T_6 := (5L_0(4/9) + 4 \cdot 0.16)/9 > 7/40 = c$ for $n \ge 6$, and at least $T_7 := (5L_0(3/7) + 4 \cdot 0.16)/9 > 0.2219$ for $n \ge 7$.
Summing over coordinates cancels the linear terms; by Proposition~\ref{prop:fourier}, this proves~(1).

For odd $n = k + 1$ ($k$ even), put $P_k := 2^{-k}\binom{k}{k/2} = p(k)$ and $S_k := \sum_{i=1}^k \varepsilon_i$.
We have $\mathbb{E}|S_k|^3 = 2k^2 P_k$ by~\eqref{eq:majorant}, and summing $(2j - k)\binom{k}{j} = k[\binom{k-1}{j-1} - \binom{k-1}{j}]$ over $j > k/2$ gives $\mathbb{E}|S_k| = k P_k$.
Conditioning on $S_k$, adding one sign increases the unnormalized cubic moment by $3\mathbb{E}|S_k| + P_k$, yielding
\[
M_n = \frac{(2n - 1)P_k}{\sqrt{n}} = h(n - 1)\frac{2n - 1}{2\sqrt{n(n - 1)}}.
\]
Set $\delta_n := n^2(h(n) - M_n) \ge 0$.
Integrating~\eqref{eq:logh} on $[n - 1, n]$ bounds $\log(h(n)/h(n-1)) \le \frac{1}{4n(n-1)}$, while $\bigl(\frac{2n-1}{2\sqrt{n(n-1)}}\bigr)^2 = 1 + \frac{1}{4n(n-1)}$.
Using $\log(1 + v) \ge v - v^2/2$, $1 - e^{-v} \le v$, and $h(n) \le \mu < 1.596$ gives
\begin{equation}\label{eq:delta}
\delta_n \;\le\; 1.596\left(\frac{n}{8(n - 1)} + \frac{1}{64(n - 1)^2}\right),
\end{equation}
which decreases in $n$ and is bounded by $e_7 = 0.234$, $e_9 = 0.225$, and $e_{11} = 0.220$ for $n = 7$, $n = 9$, and $n \ge 11$.
When $D \ge 6$ and $n \ge 7$, $(T_7 - c) \cdot 6 > 0.234 > \delta_n$, establishing~\eqref{eq:main}.

For $D \in [H(n), 6]$, every coordinate satisfies $x_i \le x(D) := (1 + \sqrt{D})/N < 1/2$ (where $N \in \{7, 9, 11\}$ is the smallest dimension in its regime), so applying~\eqref{eq:curv} along each whole tangent segment proves~\eqref{eq:main} whenever $\psi_N(D) := D[L_0(x(D)) - c] \ge e_N$.
Since $L_0' < 0$ and $L_0'' < 0$ on $[0, 1/2]$, we have $\psi_N''(D) = \frac{3}{2}L_0'(x)x' + D L_0''(x)(x')^2 < 0$ on $(0, 6]$, so $\psi_N$ achieves its minimum on $[H(N), 6]$ at an endpoint.
With $j = 100 D$ and the upward rational root bound $r = (1 + \lfloor\sqrt{j \cdot 10^{10}}\rfloor)/10^6$, exact rational evaluation (`far\_certificate` in \texttt{explore.py}) verifies $D[L_0((1 + r)/N) - 7/40] > e_N$ at all six endpoints, proving~(2).
\end{proof}

\section{Near-Flat Odd Dimensions \texorpdfstring{$n \ge 7$}{n >= 7}}

Assume $n \ge 7$ is odd and $D \le H(n)$.
Work in rescaled coordinates $u_i := \sqrt{n x_i} \ge 0$ ($\sum u_i^2 = n$).
Fix a least coordinate $d$; among the other $n - 1$ coordinates, let $A$ be a largest and $B$ a smallest.
Necessarily $B > 1/5$, since two coordinates $\le 1/5$ would give $D \ge 2(1 - 1/25)^2 = 1.8432 > 1.75$.

\begin{lemma}[Two-coordinate derivative and central-slice bound]\label{lem:slices}
When $A > B$, set $h_+ := A + B$, $l := A - B$, and for $r \ge 0$ define
\[
K(r) := \frac{(h_+ - r)_+^2}{h_+} - \frac{(l - r)_+^2}{l} \qquad (\text{second term } 0 \text{ if } l = 0).
\]
Put $m := n - 3$ (even), $z := 1/m$, let $b_1, \dots, b_m \in [B, A]$ be the remaining coordinates, and set $w := m^{-1}\sum_{i=1}^m b_i^2 = 1 + z(3 - d^2 - B^2 - A^2)$.
\begin{enumerate}[label=\textup{(\arabic*)}]
\item Letting $R$ be the residual signed sum of the $n - 2$ coordinates outside $\{A, B\}$,
\begin{equation}\label{eq:deriv}
\frac{\partial_{x_B} F - \partial_{x_A} F}{2(x_A - x_B)} = \frac{3\sqrt{n}}{8AB}\,\mathbb{E}\bigl[K(|R|)\bigr].
\end{equation}
\item $K$ is nonnegative and nonincreasing, and $v \mapsto K(\sqrt{v})$ is $C^1$ and convex on $[0, \infty)$.
\item Setting $E := H(n) - (d^2 - 1)^2 - (B^2 - 1)^2 - (A^2 - 1)^2$,
\begin{equation}\label{eq:y0y2}
\begin{aligned}
v &:= \min\!\left\{\frac{E}{(1 - z)(1 + B)^2},\; \frac{(A - B)^2}{4z(1 - z)}\right\}, \\
Y_0 &:= \sqrt{d^2 + v}, \qquad
Y_2 := \sqrt{(2\sqrt{w} - d)^2 + (1 - 4z^2)v},
\end{aligned}
\end{equation}
the right-hand side of~\eqref{eq:deriv} is at least
\begin{equation}\label{eq:fourvar}
\frac{3}{8AB}\cdot 0.797\,\sqrt{\frac{1 + 3z}{1 + z}}\left[K(Y_0) + \frac{K(Y_2)}{1 + 2z}\right].
\end{equation}
\end{enumerate}
\end{lemma}

\begin{proof}
(1) Put $H(r, t) := \frac{1}{2}(|r + t|^3 + |r - t|^3)$.
Conditional on $R$, the unnormalized two-coordinate moment is $\frac{1}{2}[H(|R|, h_+) + H(|R|, l)]$; its gradient difference in $(A^2, B^2)$ divided by $2(A^2 - B^2)$ equals $[H_t(|R|, h_+)/h_+ - H_t(|R|, l)/l]/(8AB)$.
Using $H_t(r, t)/t = 6r + 3(t - r)_+^2/t$ and rescaling by $\sqrt{n}$ yields~\eqref{eq:deriv}.

(2) On $[0, l^2]$, $K(\sqrt{v}) = h_+ - l + (1/h_+ - 1/l)v$ is affine; on $[l^2, h_+^2]$, $K(\sqrt{v}) = h_+ - 2\sqrt{v} + v/h_+$ has matching slopes at both endpoints and second derivative $\frac{1}{2}v^{-3/2} > 0$; on $[h_+^2, \infty)$ it vanishes identically.

(3) Write $\bar{b} = m^{-1}\sum b_i$ and $V = \frac{m}{m - 1}\sum (b_i - \bar{b})^2$.
Conditional on the unweighted sign sum $j = \sum_{i=1}^m \varepsilon_i$, exchangeability gives $\Var(\varepsilon_i \mid j) = 1 - j^2/m^2$ and $\Var(\sum \varepsilon_i \mid j) = 0$, determining the common off-diagonal covariance and yielding $\mathbb{E}[\sum b_i \varepsilon_i \mid j] = j\bar{b}$ and $\Var(\sum b_i \varepsilon_i \mid j) = (1 - j^2/m^2)V$.
On the slice $j = 0$ (probability $P_m$), averaging over the sign of $d$ gives $\mathbb{E}[R^2 \mid j = 0] = d^2 + V$.
On the slices $j = \pm 2$, retaining only the sign of $d$ opposite to $j$ keeps total probability $P_m/(1 + 2z)$ and conditional second moment $(2\bar{b} - d)^2 + (1 - 4z^2)V$, where $0 \le 2\bar{b} - d \le 2\sqrt{w} - d$.
Since $\sum(b_i - \bar{b})^2 \le \sum(b_i - 1)^2 \le E/(1 + B)^2$ and also $\sum(b_i - \bar{b})^2 \le m(A - B)^2/4$, we have $V \le v$.
Applying Jensen's inequality to $v \mapsto K(\sqrt{v})$ on these slices, dropping all other slices ($K \ge 0$), and using $P_m\sqrt{m + 1} \downarrow \sqrt{2/\pi} > 0.797$ (successive ratio $\sqrt{(m+1)(m+3)}/(m+2) < 1$) gives~\eqref{eq:fourvar}.
\end{proof}

\begin{theorem}[Near-flat odd $n \ge 7$]\label{thm:near}
Inequality~\eqref{eq:main} holds for every odd $n \ge 7$ whenever $D \le H(n)$.
\end{theorem}

\begin{proof}
\emph{Step~1 (Universal-$n$ interval certificate).}
Every feasible configuration satisfies $d \in [0, 1]$, $B \in [1/5, 3/2]$ (as $(n - 1)B^2 \le n \implies B \le \sqrt{7/6} < 3/2$), $A \in [1, 8/5]$ (as $A^2 \le 1 + \sqrt{1.75} < (8/5)^2$), $d \le B \le A$, $E \ge 0$, and $B^2 \le w \le A^2$.
All odd $n \ge 7$ are covered simultaneously by the three $z$-intervals in Table~\ref{tab:near}, with $z \in [0, 0.125]$ covering every $1/(n - 3)$ for odd $n \ge 11$.
The function \texttt{near\_certificate} in \texttt{explore.py} represents endpoints as integer multiples of $10^{-6}$ and evaluates outward integer-arithmetic enclosures of $E_+$, $[w_-, w_+]$, $v_+$ (omitting the second branch of~\eqref{eq:y0y2} when $z_- = 0$), $Y_0, Y_2$, and the monotone lower bound
\[
K_-(y) := \max\!\left(0,\; \frac{(A_- + B_- - y)_+^2}{A_- + B_-} - \frac{(\ell_+ - y)_+^2}{\ell_+}\right), \qquad \ell_+ := \max(0, A_+ - B_-),
\]
valid because $K$ decreases in $y$ and $(t - y)_+^2/t$ increases in $t > 0$.
Branch-and-bound bisection terminates with every box either pruned (infeasible) or certified to satisfy~\eqref{eq:fourvar} $\ge 7/40$ (Table~\ref{tab:near}).
Hence $(\partial_{x_B} F - \partial_{x_A} F)/(2(x_A - x_B)) \ge c > C$ whenever $A > B$.

\emph{Step~2 (Equalization path and endpoint comparison).}
Continuously transfer squared mass from the largest unpinned level to the smallest unpinned level, sharing rate equally within tied groups and merging levels whenever they meet.
Along this path, $\sum x_i = 1$, $q$ decreases (preserving $D \le H(n)$), $d$ stays a minimum, every unpinned coordinate stays $\ge B > 1/5 > 0$ (so $F$ is $C^2$ between mergers), and Step~1 implies $F + Cq$ is nondecreasing.
After finitely many mergers, we arrive at a vector proportional to $(1, \dots, 1, t)$ with $k = n - 1$ (even) copies of~$1$ and $t \in [0, 1]$.
By Section~3's binomial identities (only $S_k = 0$ contributes $t^3$), its moment and fourth-power mass are
\[
f(t) = P_k\frac{2k^2 + 3kt^2 + t^3}{(k + t^2)^{3/2}}, \qquad Q(t) = \frac{k + t^4}{(k + t^2)^2}.
\]
On $(0, 1)$, $Q'(t) < 0$ and
\[
\frac{f'(t)}{-Q'(t)} = \frac{3P_k}{4}\,\frac{t\sqrt{k + t^2}}{1 + t}
\]
is strictly increasing in $t$, so $f(t) + CQ(t)$ attains its maximum on $[0, 1]$ at $t = 0$ or $t = 1$.
The secant quotient comparing $t = 1$ with $t = 0$ is
\[
k(k + 1)(M_{k+1} - M_k) = P_k\sqrt{k}\,\frac{r}{(1 + r)^2}, \qquad r := \sqrt{1 + 1/k}.
\]
As even $k \ge 2$ increases, $P_k\sqrt{k}$ strictly increases (successive ratio $(k + 1)/\sqrt{k(k + 2)} > 1$) and $r$ decreases toward~$1$, strictly increasing $r/(1 + r)^2$.
At $k = 2$ this quotient equals $C$, so $f(0) + CQ(0) \le f(1) + CQ(1) = M_n + C/n$, completing the proof.
\end{proof}

\begin{table}[ht]
\centering
\begin{tabular}{@{}lrrr@{}}
\toprule
Odd dimensions covered & Threshold $H$ & Interval for $z = 1/(n - 3)$ & Visited boxes \\
\midrule
$n = 7$ & $1.75$ & $[0.250000,\, 0.250000]$ & $435$ \\
$n = 9$ & $1.45$ & $[0.166666,\, 0.166667]$ & $427$ \\
All odd $n \ge 11$ & $1.24$ & $[0.000000,\, 0.125000]$ & $5{,}747$ \\
\bottomrule
\end{tabular}
\caption{Exact integer-arithmetic branch-and-bound certificates for~\eqref{eq:fourvar} $\ge 7/40$.}\label{tab:near}
\end{table}

\section{Exact Cone Certificates for Dimensions \texorpdfstring{$1 \le n \le 5$}{1 <= n <= 5}}\label{sec:small}

For $n = 1$, both sides of~\eqref{eq:main} are zero.
For $2 \le n \le 5$, let $a_1 \ge \cdots \ge a_n \ge 0$ be unnormalized, put $N = \sum a_i^2$, $Q = \sum a_i^4$, $b = 2^{n-1}$, $\mathcal{S} = \{1\} \times \{\pm 1\}^{n-1}$, and $h_0 = \sum_{\sigma \in \mathcal{S}} |\sum_i \sigma_i|^3$, so $M_n = h_0/(bn\sqrt{n})$.
For each row $(n, k)$ of Table~\ref{tab:cones}, write $\sqrt{n}\,C_{n,k} = (A + B\sqrt{nk})/d$ with integers $A, B, d$ ($d > 0$).
The sharp dimension-$n$ inequality $M_n - F(x) \ge C_{n,k}\Delta_n(a)$ is equivalent to
\begin{equation}\label{eq:cone-ineq}
\sqrt{n}\,nd\sqrt{N}\sum_{\sigma \in \mathcal{S}} |\sigma \cdot a|^3
\;\le\; d h_0 N^2 + (A + B\sqrt{nk})\,b(N^2 - nQ),
\end{equation}
whose right-hand side is strictly positive for $N > 0$ (as $M_n \ge 1$ by Jensen and $0 < C_{n,k} < 1$).

On a simplicial cone $a = \sum_{j=1}^n t_j r_j$ ($t_j \ge 0$), split each linear form $\sigma \cdot a = P_\sigma(t) - N_\sigma(t)$ into its nonnegative- and negative-coefficient parts in $t$, and set $U := \sum_{\sigma \in \mathcal{S}} (P_\sigma^3 + N_\sigma^3) \ge \sum_{\sigma \in \mathcal{S}} |\sigma \cdot a|^3$.
Squaring~\eqref{eq:cone-ineq} with $U$ on the left yields the polynomial
\[
\bigl[d h_0 N^2 + (A + B\sqrt{nk})\,b(N^2 - nQ)\bigr]^2 - d^2 n^3 N U^2 \;\in\; \mathbb{Z}[\sqrt{nk}][t_1, \dots, t_n].
\]
If every monomial coefficient $u + v\sqrt{nk}$ is nonnegative (verified by exact sign-preserving integer squaring), then~\eqref{eq:cone-ineq} holds throughout the cone.
Starting from the ordered cone $\{r_j = (1^j, 0^{n-j}) : 1 \le j \le n\}$ and recursively splitting along the first mixed-sign hyperplane $\sigma \cdot a = 0$ (inserting $(q r_j + p r_l)/\gcd$ where $\sigma \cdot r_j = -p < 0 < q = \sigma \cdot r_l$, or bisecting the longest Euclidean edge on $\sum a_i = 1$ if no mixed sign exists), \texttt{chamber\_certificate} in \texttt{explore.py} terminates with the exact leaf counts in Table~\ref{tab:cones}.
Since $C_{2,1}, C_{4,2}, C_{5,4} > 7/40 > C$ and $C_{3,2} = C$ (Section~1), combining Table~\ref{tab:cones} with Theorems~\ref{thm:far} and~\ref{thm:near} completes the proof of Theorem~\ref{thm:main}.

\begin{table}[ht]
\centering
\begin{tabular}{@{}rrrc@{}}
\toprule
$n$ & Support $k$ & Optimal constant $C_{n,k} = (M_n - M_k)/(1/k - 1/n)$ & Certified cones \\
\midrule
$2$ & $1$ & $2(\sqrt{2} - 1) \approx 0.828427$ & $1$ \\
$3$ & $2$ & $5\sqrt{3} - 6\sqrt{2} \approx 0.174973$ & $2$ \\
$4$ & $2$ & $6 - 4\sqrt{2} \approx 0.343146$ & $3$ \\
$5$ & $4$ & $27\sqrt{5}/2 - 30 \approx 0.186918$ & $41$ \\
\bottomrule
\end{tabular}
\caption{Exact integer-radical polynomial cone certificates for $2 \le n \le 5$.}\label{tab:cones}
\end{table}

\section{Reproducible Verification}

Running \texttt{explore.py} in \texttt{solutions/rademacher\_fourth\_moment/} via \texttt{.venv/bin/python -B} from the repository root deterministically verifies all $47$ polynomial cones, the $6$ exact rational far-region endpoints, and all $6{,}609$ integer-arithmetic interval boxes.
Independently of the proof, it also executes $32$ seeded BFGS sphere searches per dimension for $2 \le n \le 13$ with exact integer re-evaluation at $75$-digit precision, recovering $0.174972663606$ at $n = 3$ and minimizing supports $4, 6, 8, 10, 12$ for odd $n = 5, 7, 9, 11, 13$. All assertions pass.

\begin{thebibliography}{9}

\bibitem{GaoQian2026}
P.~Gao and J.~Qian.
\emph{Fourth-Moment Geometry of Rademacher Sums}.
Preprint \href{https://arxiv.org/abs/2608.17802}{arXiv:2608.17802}, August 2026.

\end{thebibliography}

\end{document}
